Item availability in a restaurant point-of-sale (POS) backend must reflect orders that deplete shared ingredients. Redis cache-aside can reduce database query activity, while expiration and deletion expose different freshness risks. This study compares PostgreSQL-only reads, TTL-only caching, deletion after commit, and deletion with generation-checked refills. In a controlled schedule, a delayed SQL snapshot made all ten follow-up reads stale with plain deletion. A generation check blocked that refill in all ten corresponding trials. In a separate last-portion workload, TTL-only caching prompted 80 checkout attempts across 80 actual browsing sessions after inventory had sold out. PostgreSQL rejected all 80 attempts. At 4{,}800 offered requests/s in ten repeated k6 trials per mode, TTL-only caching reduced the median per-trial item GET P95 from 28.67 to 4.14\,ms, while order POST P95 was 54.30\,ms compared with 35.31\,ms in PostgreSQL-only mode. These measurements describe one Linux host under actual trial workloads. The source code, recorded data, and analysis support reanalysis.
